Dans le BAC S de Juin 2018 (Métropole), la partie A de l'exercice de Spécialité Mathématiques s'intéressait à l'équation diophantienne \ED\ sur $\NN^2$ qui admet la solution évidente $(x \,; y) = (3 \,; 1)$. L'exerice introduit une matrice \squote{magique}
$A =
\begin{pmatrix} 
  3 & 8  \\ 
  1 & 3 
\end{pmatrix}$
permettant la construction de nouvelles solutions $(x_n \,; y_n)$ de façon récursive via
$\begin{pmatrix} 
  x_{n+1} \\ 
  y_{n+1} 
\end{pmatrix}
=
A
\begin{pmatrix} 
  x_{n} \\ 
  y_{n} 
\end{pmatrix}
$.
Très bien ! Mais, comment peut-on découvrir la matrice \squote{magique} $A$ ?